\section{Relativistic Material Collisions in Spherical Spacetime}
\hypertarget{sec:relativistic-material-collisions}{}
\label{sec:relativistic-material-collisions}

The relativistic collision library adds elastic material-particle scattering
to the Schwarzschild-coordinate (SC) and advanced Eddington--Finkelstein (EF)
representations developed above. A material particle has rest mass
\(m_\chi>0\). A statistical packet carries represented rest mass \(W\), so
it represents \(q=W/m_\chi\) particles. Its momentum and its statistical
weight have separate physical roles. Local scattering changes momentum,
while conservative packet splitting preserves represented particle number
and rest mass.

The library first evaluates collision rates on the chosen spacetime slice,
then scatters each accepted pair in a local orthonormal frame. Persistent cell
clocks carry coordinate time, remaining proposal waiting depth, and random
state between coupled transport and collision steps. The constant-cross-section spherical sampler
and local amplitude interface have distinct entry points. The derivation
below establishes their measures and normalizations before introducing those
source calls.

\subsection{Invariant equation and represented mass}

Use signature \((-+++)\), geometrized units \(G=c=1\), and physical
four-momentum \(k^\alpha=m_\chi u^\alpha\), where
\(u_\alpha u^\alpha=-1\). Let \(f_N(x,k)\) be the scalar number
distribution in physical momentum. In a local orthonormal frame,
\(E=k^{\hat0}>0\), and
\[
 \dd\Pi=\frac{\dd^3 k}{E},\qquad
 J^\alpha=\int k^\alpha f_N\dd\Pi,\qquad
 T^{\alpha\beta}=\int k^\alpha k^\beta f_N\dd\Pi.
 \tag{REB1}
\]
The invariant particle count crossing an oriented hypersurface is
\(\dd N=f_N k^\alpha\dd\Sigma_\alpha\dd\Pi\), with the surface
orientation chosen to make the future crossing count positive.
For a spacelike slice this is equivalently
\(\dd N=f_N(-k\cdot n)\dd\Sigma\dd\Pi\), where \(n\) is its
future unit normal. These expressions supply the different SC and EF
counting measures.

For classical elastic \(2\to2\) scattering at one spacetime event, the
Boltzmann equation is
\[
 \begin{split}
 k^\alpha\partial_\alpha f_N
 -\Gamma^i{}_{\alpha\beta}k^\alpha k^\beta
       \frac{\partial f_N}{\partial k^i}
 &=C_N[f_N],\\
 C_N[f_{N1}]&=\int\dd\Pi_2\,\mathcal F
   \int\dd\Omega\,\frac{\dd\Sigma}{\dd\Omega}
       (f_{N3}f_{N4}-f_{N1}f_{N2}),\\
 \mathcal F&=\sqrt{(k_1\cdot k_2)^2-m_\chi^4}.
 \end{split}\tag{REB2}
\]
Here \(\Gamma^i{}_{\alpha\beta}\) is the coordinate connection, and
\(\dd\Sigma/\dd\Omega\) is the physical differential cross section. The solid angle \(\dd\Omega\)
specifies the outgoing relative momentum in the pair centre-of-momentum frame.
The connection governs transport between events. At a collision event, a
local tetrad supplies an orthonormal observer frame in which the two incoming
momenta are scattered. The collision integral uses the metric at that same
event. The supplied
interaction obeys inverse-scattering reciprocity, which gives the stated
gain-minus-loss form. The spherical sampler uses
\(\dd\Sigma/\dd\Omega=\Sigma/(4\pi)\), with constant material area
\(\Sigma\) in code units. The M\o ller flux factor
\(\mathcal F/(E_1E_2)\) can exceed one. Each particle velocity still
has magnitude below one.

The preceding EV chapters use specific momentum \(u=k/m_\chi\) and a
represented-rest-mass distribution \(f\). Changing momentum variables
gives
\[
 f=m_\chi^4 f_N,\qquad
 \dd\Pi=m_\chi^2\dd\Pi_u,\qquad
 \mathcal F=m_\chi^2\mathcal F_u,\qquad
 \mathcal F_u=\sqrt{(u_1\cdot u_2)^2-1}.
 \tag{REB3}
\]
Consequently \(L_u[f]=m_\chi^{-1}C_u[f]\), where \(L_u\) is the
transport operator on the left of (REB2), expressed in specific momentum and \(C_u\) has the form in (REB2)
with \(\dd\Pi_u,\mathcal F_u,f\). The material coefficient in this
equation is therefore \(\Sigma/m_\chi\). The implemented stress retains
the same rest-mass packet weights as the EV solver.

\subsection{SC proper volume and collision rate}

Schwarzschild coordinate time is \(\tau\). A stationary local observer
measures proper interval \(e^\mu\dd\tau\). In the SC metric, the local
specific radial momentum is \(w\), and
\(\eta=L^2/2\) uses specific angular momentum \(L\). Define
\(\epsilon=\sqrt{1+w^2+2\eta/r^2}\). The slice normal measures energy
\(E=m_\chi\epsilon\), and \(u^\tau=e^{-\mu}\epsilon\).
Dividing the invariant equation by \(u^\tau\) gives
\[
 \begin{split}
 \partial_\tau f
 +e^{\mu-\lambda}\frac{w}{\epsilon}\partial_r f
 -\left[\lambda_\tau w+e^{\mu-\lambda}\mu_r\epsilon
       -e^{\mu-\lambda}\frac{2\eta}{r^3\epsilon}\right]
       \partial_w f
 =\frac{e^\mu}{m_\chi\epsilon}C_u[f].
 \end{split}\tag{REB4}
\]
Writing \(\beta\) for tangential momentum azimuth gives
\(\dd\Pi_u=\dd w\dd\eta\dd\beta/(r^2\epsilon)\).
Integration over position angles and \(\beta\) yields represented rest
mass \(\dd M_{\rm rest}=8\pi^2e^\lambda f\dd r\dd w\dd\eta\).

In a homogeneous proper cell of volume \(V\), the empirical number
measure is \(f_N\dd\Pi=\sum_i q_i\delta_i/(VE_i)\). Multiplication
by proper four-volume \(e^\mu V\dd\tau\) then gives the unordered
physical pair rate \(q_iq_j e^\mu\Sigma\mathcal F/(VE_iE_j)\).
An accepted packet event scatters \(\min(q_i,q_j)\) particles from each
packet. Its event hazard per SC coordinate time is thus
\[
 \Gamma_{ij}^{\rm SC}
 =\max(W_i,W_j)\frac{\Sigma}{m_\chi}
   \frac{e^\mu\mathcal F_u}{V\epsilon_i\epsilon_j}.
 \tag{REB5}
\]
The finite radial-cell adapter computes \(V\) and the volume-averaged
lapse from the reconstructed SC metric. The reported metric ranges sample
its eight quadrature points. This is a homogeneous-cell approximation.
Its empty-Minkowski limit gives
\(V=4\pi(r_{\rm R}^3-r_{\rm L}^3)/3\) and mean lapse one.

\subsection{EF null flux and collision rate}

For \(\dd s^2=-ab^2\dd v^2+2b\dd v\dd r+r^2\dd\Omega^2\), retain
the preceding EF chapter's positive specific canonical momentum \(P=u_r\).
Define \(\chi=1+2\eta/r^2\). The future mass shell and the prescribed
radial observer are
\[
 \begin{gathered}
 u_v=-\frac b2\left(aP+\frac\chi P\right),\qquad
 u^v=\frac P b,\qquad
 u^r=\frac12\left(aP-\frac\chi P\right),\\
 U=\frac1b\partial_v+\frac{a-1}{2}\partial_r,\qquad
 e_{\hat r}=\frac1b\partial_v+\frac{a+1}{2}\partial_r.
 \end{gathered}\tag{REB6}
\]
For finite \(a\) and \(b>0\), this observer is unit timelike and its
radial vector is unit spacelike, including at \(a=0\). It measures
\[
 \epsilon_U=\frac12\left(P+\frac\chi P\right),\qquad
 w_U=\frac12\left(P-\frac\chi P\right),\qquad
 \boldsymbol{k}_\perp=
   \frac{m_\chi\sqrt{2\eta}}r(\cos\beta,\sin\beta).
 \tag{REB7}
\]
The physical energy and radial momentum are \(m_\chi\epsilon_U\) and
\(m_\chi w_U\). Projecting an outgoing momentum gives
\(P=(E+k^{\hat r})/m_\chi\). The code evaluates the equivalent
\(P=\chi/(\epsilon_U-w_U)\) for negative \(w_U\) to avoid cancellation.

The Jacobian \(\dd w_U/\dd P=\epsilon_U/P\) determines both the
shell measure and the null-slice count:
\[
 \dd\Pi_u=\frac{\dd P\dd\eta\dd\beta}{r^2P},\qquad
 \dd M_{\rm rest}=8\pi^2 f\dd r\dd P\dd\eta.
 \tag{REB8}
\]
The directed null surface has count
\(\dd N=J^v b r^2\dd r\dd\Omega\). With coordinate shell volume
\(V_0=4\pi r^2\dd r\), a packet contributes
\(J^v=q/(bV_0)\). Its empirical invariant momentum measure is therefore
\(f_N\dd\Pi=\sum_iq_i\delta_i/(bV_0k_i^v)\). Combining this measure
with invariant four-volume \(bV_0\dd v\) gives
\[
 \Gamma_{ij}^{\rm EF}
 =\max(W_i,W_j)\frac{\Sigma}{m_\chi}
   \frac{b\mathcal F_u}{V_0P_iP_j}.
 \tag{REB9}
\]
The EF rate is therefore normalized by the particles crossing a null slice.
The SC rate in (REB5) instead counts particles within a proper spatial cell.
These distinct counts explain the different coordinate-time prefactors.

The coordinate Boltzmann equation follows by dividing by \(u^v=P/b\):
\[
 \begin{split}
 \partial_v f+\frac b2\left(a-\frac\chi{P^2}\right)\partial_r f
 +\left[-\frac P2(ab_r+ba_r)-\frac{b_r\chi}{2P}
           +\frac{2b\eta}{Pr^3}\right]\partial_P f
 =\frac b{m_\chi P}C_u[f].
 \end{split}\tag{REB10}
\]
The same null measure gives the implemented constraint sources
\[
 T_{rr}=\frac1{V_0}\sum_iW_iP_i,\qquad
 S=\frac1{V_0}\sum_iW_i\frac{\chi_i}{P_i}.
 \tag{REB11}
\]
At one radius, elastic four-momentum conservation preserves both sums
instantaneously. The changed momentum distribution subsequently changes
transport and the other stress components. Passive-coordinate tests compare
the same physical event, observer and surface. SC and EF time slices are
different physical surfaces, so their coordinate-time hazards describe
their respective crossing ensembles.

\subsection{Coupled inputs, streaming and cell geometry}

The entry points are \code{eb::advance_sc} and \code{eb::advance_ef} in
\path{cpu/include/ev/relativistic_boltzmann/coupled_evolution.hpp}.
Their state contains the marker ensemble, global coordinate time, one
persistent collision clock per radial cell, and the next unused packet
identifier. The caller supplies the field grid, EV evolution controls and
\code{eb::SplitControls}. The latter owns material mass, constant code
area, packet/proposal limits and the finite-cell policy.

The maintained state definition makes this ownership explicit:
\par\medskip\noindent
\begin{minipage}{\linewidth}
\codesource{Coupled material state}{\path{cpu/include/ev/relativistic_boltzmann/coupled_evolution.hpp}, lines 28--36}{Carry the marker ensemble and all persistent clocks on the selected slice}
\begin{lstlisting}[language=C++,showstringspaces=false]
template<class Marker> struct CoupledState {
    std::vector<Marker> markers;
    std::vector<Cell> collision_clocks;
    double time = 0;
    std::uint64_t next_id = 1;
};

using SCState = CoupledState<ev::WeightedPICMarker>;
using EFState = CoupledState<ev::eddington_finkelstein::Marker>;
\end{lstlisting}
\end{minipage}

A caller supplies its physical distribution, material mass and area in
\code{SplitControls::model}, and sizes \code{collision_clocks} to the field-cell
count. Each cell starts at the global coordinate time with a recorded random
state, and \code{next_id} exceeds every input identifier. The default local
policy uses a co-located population. Setting
\code{allow_radial_cell_approximation=true} declares the finite homogeneous
radial-cell approximation. The caller then invokes
\code{eb::advance_sc} with the state, requested interval, grid, evolution
options and split controls, or its EF counterpart.
Both require the isolated vacuum outer-boundary guard. EF initial data are
constructed on the chosen null surface using the transformation derived in
the preceding chapter.

Each entry point first copies the caller's state and advances the
maintained self-consistent EV solver. Its returned \code{accepted_dt}
is the collision exposure, including any shortening by a transport retry.
Thus SC exposure is a Schwarzschild coordinate interval and EF exposure
is an advanced-time interval. The endpoint field remains fixed throughout
the collision suboperator. This streaming--collision composition is
first order. A zero area or zero accepted interval returns the direct EV
update while advancing the persistent coordinate clocks consistently.

The SC streaming call is:
\par\medskip\noindent
\begin{minipage}{\linewidth}
\codesource{Self-consistent SC streaming}{\path{cpu/include/ev/relativistic_boltzmann/coupled_evolution.hpp}, lines 108--115}{Use the accepted EV update and retain the direct collision-off control}
\begin{lstlisting}[language=C++,showstringspaces=false]
    out.streaming = ev::advance_weighted_einstein_vlasov_pic_self_consistent(s.markers,dt,grid,evolution);
    out.field = out.streaming.endpoint_field;
    if (c.model.cross_section == 0 || out.streaming.accepted_dt == 0) {
        s.time += out.streaming.accepted_dt;
        for (auto& clock : s.collision_clocks) clock.coordinate_time = s.time;
        caller = std::move(s);
        return out;
    }
\end{lstlisting}
\end{minipage}

The EF entry point obtains its accepted interval from the maintained null
streaming update:
\par\medskip\noindent
\begin{minipage}{\linewidth}
\codesource{Self-consistent EF streaming}{\path{cpu/include/ev/relativistic_boltzmann/coupled_evolution.hpp}, lines 185--192}{Use the accepted advanced-time interval and preserve the direct collision-off branch}
\begin{lstlisting}[language=C++,showstringspaces=false]
    out.streaming = ef::advance_self_consistent(s.markers,dv,grid,evolution);
    out.field = out.streaming.endpoint_field;
    if (c.model.cross_section == 0 || out.streaming.accepted_dt == 0) {
        s.time += out.streaming.accepted_dt;
        for (auto& clock : s.collision_clocks) clock.coordinate_time = s.time;
        caller = std::move(s);
        return out;
    }
\end{lstlisting}
\end{minipage}

The streamed radius determines current cell membership through a binary
search of increasing field faces. The method fills each containing cell's packet list directly. Marker identifiers and source labels retain their own
meaning. This assignment supplies local collision partners, while WPIC
gravity continues to use deposited stress on the field mesh. The incoming
SC material markers must be active with zero unresolved subhalo internal
energy and dispersion. EF retains every active material marker.
\par\medskip\noindent
\begin{minipage}{\linewidth}
\codesource{Current-radius cell assignment}{\path{cpu/include/ev/relativistic_boltzmann/coupled_evolution.hpp}, lines 57--62}{Find the containing radial cell after the accepted drift}
\begin{lstlisting}[language=C++,showstringspaces=false]
inline std::size_t radial_cell(double r, const ev::WeightedPICFieldGrid& grid) {
    const auto upper = std::upper_bound(grid.mesh.edge.begin(),grid.mesh.edge.end(),r);
    rb::require(upper != grid.mesh.edge.begin() && r < grid.mesh.edge.back(),
                "Matter reached the isolated collision-domain boundary.");
    return static_cast<std::size_t>(upper-grid.mesh.edge.begin()-1);
}
\end{lstlisting}
\end{minipage}

Before refilling each cell's packets, the code checks that its coordinate
clock equals the global time. Clearing membership preserves its random
state, residual optical depth, carried number-weight bound and counters:
\par\medskip\noindent
\begin{minipage}{\linewidth}
\codesource{Persistent clock ownership}{\path{cpu/include/ev/relativistic_boltzmann/coupled_evolution.hpp}, lines 84--93}{Refill membership while retaining sampler state}
\begin{lstlisting}[language=C++,showstringspaces=false]
inline void prepare_clocks(CoupledState<Marker>& s, const ev::WeightedPICFieldGrid& grid) {
    rb::require(s.collision_clocks.size() == grid.radius.size(),"One persistent collision clock per cell is required.");
    for (auto& clock : s.collision_clocks) {
        rb::require(clock.coordinate_time == s.time,"The material and collision clocks disagree.");
        // Refill packet membership from the transported population below.
        // Random state, residual optical depth and the carried majorant persist.
        clock.packets.clear();
    }
}

\end{lstlisting}
\end{minipage}

An entering packet raises the cell's carried \(q_{\max}\) when needed.
For SC, eight-point radial Gauss--Legendre quadrature of the endpoint
metric supplies proper volume \(V=\int4\pi r^2e^\lambda\dd r\) and
mean lapse \(V^{-1}\int e^\mu\dd V\). Its sampled ranges describe those
quadrature nodes. Each occupied cell then uses proper volume and mean lapse from that
endpoint field. The following call advances its persistent collision clock
by the returned accepted interval.
\par\medskip\noindent
\begin{minipage}{\linewidth}
\codesource{SC cell exposure}{\path{cpu/include/ev/relativistic_boltzmann/coupled_evolution.hpp}, lines 138--146}{Use endpoint geometry and the actual accepted coordinate interval}
\begin{lstlisting}[language=C++,showstringspaces=false]
        const auto geometry = rb::sc_cell_geometry(grid,out.streaming.endpoint_field,i);
        const auto g = collision_geometry(Slice::schwarzschild,geometry.proper_volume,
            geometry.mean_lapse,geometry.outer_radius-geometry.inner_radius,c);
        const auto before = cell.collisions;
        const auto nonlocal = cell.radially_nonlocal_collisions;
        advance_cell(cell,out.streaming.accepted_dt,c.model,g);
        out.collisions += cell.collisions-before;
        out.nonlocal_collisions += cell.radially_nonlocal_collisions-nonlocal;
        s.next_id = cell.next_id;
\end{lstlisting}
\end{minipage}



For EF, the code uses the exact coordinate shell volume
\(V_0=4\pi(r_{i+1}^3-r_i^3)/3\) and the endpoint value of \(b\) at
that cell's centre. The null momentum factors supply the remaining flux
conversion in (REB9):
\par\medskip\noindent
\begin{minipage}{\linewidth}
\codesource{EF endpoint geometry and exposure}{\path{cpu/include/ev/relativistic_boltzmann/coupled_evolution.hpp}, lines 211--217}{Use coordinate volume and the endpoint centre value of b for the accepted interval}
\begin{lstlisting}[language=C++,showstringspaces=false]
        const double b = out.streaming.endpoint_field.metric.b[i];
        const double V0 = grid.mesh.shell_volume[i];
        const auto g = collision_geometry(Slice::advanced_null,V0,b,
            grid.mesh.edge[i+1]-grid.mesh.edge[i],c);
        const auto before = cell.collisions;
        const auto nonlocal = cell.radially_nonlocal_collisions;
        advance_cell(cell,out.streaming.accepted_dt,c.model,g);
\end{lstlisting}
\end{minipage}

\subsection{Proposal clock, conservative update and source feedback}

Let \(N\) be the current packet count and \(q_{\max}\) a carried upper
bound on packet number weight. Splitting never increases a packet's
weight. On SC set \(D_{\max}=1\). On EF recompute
\(D_{\max}=\max_i(\epsilon_{Ui}/P_i)\) after every accepted state
change. The implementation evaluates this angle-independent bound at
the start of every proposal loop. The inequality
\(\mathcal F_u/(P_iP_j)\leq2D_{\max}^2\) provides a valid EF
envelope. The current momenta therefore determine the EF rate bound.
With \((A,V)=(e^\mu,V)\) on SC and \((b,V_0)\) on EF, the proposal
rate and uniform-unordered-pair acceptance are
\[
 \begin{split}
 R&=\frac{A\Sigma}{V}q_{\max}N(N-1)D_{\max}^2,\\
 \mathcal P_{ij}&=\frac{\max(q_i,q_j)}{q_{\max}}
       \frac{B_{ij}}{2D_{\max}^2},\qquad
 B_{ij}=\begin{cases}
  \mathcal F_u/(\epsilon_i\epsilon_j),&\mathrm{SC},\\
  \mathcal F_u/(P_iP_j),&\mathrm{EF}.
 \end{cases}
 \end{split}\tag{REB12}
\]
Multiplying by the unordered selection probability \(2/[N(N-1)]\)
recovers (REB5) or (REB9). The factor for unordered incoming pairs is
therefore already accounted for. The persistent optical depth starts at
\(-\ln U\), where \(U\) is uniform on \((0,1)\). A waiting interval
consumes \(R\Delta t\). Residual optical depth and random-generator
state persist across calls.

Write \(\mathcal T\) for carried optical depth and \(t\) for the selected
SC or EF coordinate clock. For an unfinished interval of length \(h\), the
discrete clock update is
\[
 (t,\mathcal T)\longmapsto
 \begin{cases}
 (t+h,\mathcal T-Rh),&\mathcal T>Rh,\\
 (t+\mathcal T/R,-\ln U'),&\mathcal T\leq Rh.
 \end{cases}\tag{REB12a}
\]
The second branch processes one proposal before drawing its next optical
depth \(-\ln U'\). An accepted collision can change the next EF envelope,
so the loop recomputes its rate before spending the remaining exposure.
\par\medskip\noindent
\begin{minipage}{\linewidth}
\codesource{Proposal rate and carried optical depth}{\path{cpu/include/ev/relativistic_boltzmann/spherical_collisions.hpp}, lines 178--192}{Consume only the available coordinate-time exposure}
\begin{lstlisting}[language=C++,showstringspaces=false]
        const double n = static_cast<double>(s.packets.size());
        // Unordered pair count n(n-1)/2 and the flux envelope 2 Dmax^2 produce
        // the n(n-1) factor. The time factor is the SC lapse or EF b, and volume
        // is respectively proper spatial volume or null-slice coordinate volume.
        const double rate = g.time_factor*model.cross_section/g.volume
            *s.number_majorant*n*(n-1)*Dmax*Dmax;
        rb::require(rate > 0 && std::isfinite(rate),"Unrepresentable slice collision majorant.");
        if (s.optical_depth < 0) s.optical_depth = -std::log(s.random.uniform());
        const double wait = s.optical_depth/rate;
        // Spend only the optical depth covered by this frozen interval. The
        // remaining depth survives transport and cell rebinding at the next step.
        if (wait > end-s.coordinate_time) {
            s.optical_depth -= rate*(end-s.coordinate_time);
            break;
        }
\end{lstlisting}
\end{minipage}

The next excerpt chooses two distinct packets, samples their tangential
azimuths and applies the slice factor from (REB12). The integer sampler uses rejection to keep each index uniform, then maps
the second index around the first. Two independent tangential azimuths
lift the selected spherical momenta to physical local four-momenta. A
fresh independent uniform draw tests acceptance. These draws produce
uniform unordered pairs for the identical material interaction.
\par\medskip\noindent
\begin{minipage}{\linewidth}
\codesource{Pair proposal and acceptance}{\path{cpu/include/ev/relativistic_boltzmann/spherical_collisions.hpp}, lines 197--216}{Reject a flux-envelope violation before changing packet momenta}
\begin{lstlisting}[language=C++,showstringspaces=false]
        const std::size_t i = rb::random_index(s.random,s.packets.size());
        std::size_t j = rb::random_index(s.random,s.packets.size()-1);
        if (j >= i) ++j;
        const Packet a = s.packets[i];
        const Packet b = s.packets[j];
        // Independent tangential azimuths sample the spherical distribution at
        // the representative event. Each lifted momentum still has material
        // mass m_chi, regardless of the packet rest weight.
        const rb::Four ka = lift(a,model,g.slice,2*rb::pi*s.random.uniform());
        const rb::Four kb = lift(b,model,g.slice,2*rb::pi*s.random.uniform());
        const double factor = pair_factor(a,b,ka,kb,model,g);
        s.maximum_null_flux_factor = std::max(s.maximum_null_flux_factor,
            g.slice == Slice::advanced_null ? factor : 0);
        // The number-weight ratio and exact slice flux remove the conservative
        // proposal envelope. The bound must hold before the random accept/reject
        // draw; an exceeded envelope stops the transaction.
        const double probability = std::max(a.rest_weight,b.rest_weight)/(model.material_mass*s.number_majorant)
            *factor/(2*Dmax*Dmax);
        rb::require(probability >= 0 && probability <= 1,"Slice flux envelope was exceeded.");
        if (s.random.uniform() < probability) {
\end{lstlisting}
\end{minipage}

\par\medskip\noindent
\begin{minipage}{\linewidth}
\codesource{SC and EF pair flux}{\path{cpu/include/ev/relativistic_boltzmann/spherical_collisions.hpp}, lines 146--151}{Convert the same local invariant flux to the selected slice rate}
\begin{lstlisting}[language=C++,showstringspaces=false]
inline double pair_factor(const Packet& a, const Packet& b, rb::Four ka, rb::Four kb,
                          const Model& model, const Geometry& g) {
    const double F = rb::invariant_flux(ka,kb);
    return g.slice == Slice::schwarzschild ? F/(ka.t*kb.t)
        : F/(model.material_mass*model.material_mass*a.radial*b.radial);
}
\end{lstlisting}
\end{minipage}

For an accepted event, the code forms the total physical momentum
\(K^\alpha=k_i^\alpha+k_j^\alpha\), its invariant
\(s=-K^\alpha K_\alpha\), and its centre-of-momentum velocity
\(\boldsymbol V=\boldsymbol K/K^0\). Two independent uniform draws give
\(\cos\vartheta=2U-1\) and \(\varphi=2\pi U'\), defining the isotropic
unit vector \(\boldsymbol n_*\). Set
\(\gamma_V=(1-|\boldsymbol V|^2)^{-1/2}\). Transforming the spacelike
vector \((0,\boldsymbol n_*)\) to the original local frame gives
\[
 n^0=\gamma_V\boldsymbol V\cdot\boldsymbol n_*,\qquad
 \boldsymbol n=\boldsymbol n_*+
 \frac{\gamma_V-1}{|\boldsymbol V|^2}
   (\boldsymbol V\cdot\boldsymbol n_*)\boldsymbol V.
 \tag{REB12b}
\]
At \(\boldsymbol V=0\), the continuous limit is
\(n^\alpha=(0,\boldsymbol n_*)\). The unit total momentum is
\(U_K^\alpha=K^\alpha/\sqrt s\), and the centre-of-momentum magnitude is
\(q_*=\mathcal F/\sqrt s=\sqrt{s/4-m_\chi^2}\). The code constructs
\[
 k_3^\alpha=\frac{\sqrt s}{2}U_K^\alpha+q_*n^\alpha,\qquad
 k_4^\alpha=\frac{\sqrt s}{2}U_K^\alpha-q_*n^\alpha.
 \tag{REB12c}
\]
The relations \(n\cdot U_K=0\), \(n\cdot n=1\), and
\(U_K\cdot U_K=-1\) give \(k_3^2=k_4^2=-m_\chi^2\) and
\(k_3+k_4=K\). The implementation evaluates the outgoing energy from
its spatial momentum and checks the floating-point projection against
this covariant construction. These operations are in
\path{relativistic_collision.hpp}, \code{pair_geometry},
\code{scatter_isotropic} and \code{scatter_with_direction}. It scatters represented weight \(W_{\min}\) in each packet
and retains a larger incoming packet's unscattered remainder. This update
preserves represented number and local four-momentum for unequal weights.
Packet-cap, flux-envelope and conditioning failures reject the transaction
before replacing the caller's state.

\par\medskip\noindent
\begin{minipage}{\linewidth}
\codesource{Centre-of-momentum outgoing pair}{\path{cpu/include/ev/relativistic_boltzmann/relativistic_collision.hpp}, lines 184--187}{Form both physical outgoing momenta before the guarded mass-shell projection}
\begin{lstlisting}[language=C++,showstringspaces=false]
    const Four first = g.energy1_cm*g.unit_total+g.momentum_cm*direction;
    const Four second = g.energy2_cm*g.unit_total-g.momentum_cm*direction;
    const Four first_shell = on_shell(m1,first.p);
    const Four second_shell = on_shell(m2,second.p);
\end{lstlisting}
\end{minipage}


For an accepted event, the represented-mass update is
\[
 W_i'=W_j'=\min(W_i,W_j),\qquad
 W_{\rm rem}=|W_i-W_j|.
 \tag{REB12d}
\]
The remainder carries the larger incoming packet's momentum. The two
scattered packets carry the outgoing pair momenta, so their represented
four-momentum plus the remainder equals the incoming represented total.
\par\medskip\noindent
\begin{minipage}{\linewidth}
\codesource{Elastic event and conservative weight split}{\path{cpu/include/ev/relativistic_boltzmann/spherical_collisions.hpp}, lines 219--234}{Scatter the common represented weight and preserve the untouched remainder}
\begin{lstlisting}[language=C++,showstringspaces=false]
            const auto out = rb::scatter_isotropic(model.material_mass,ka,model.material_mass,kb,s.random);
            // Scatter the equal rest weight min(W1,W2) from both packets. A heavier
            // packet keeps a positive remainder with the incoming momentum and a new
            // ID. This preserves represented rest mass and weighted local four-momentum.
            const double scattered_weight = std::min(a.rest_weight,b.rest_weight);
            s.packets[i].rest_weight = scattered_weight;
            s.packets[j].rest_weight = scattered_weight;
            project(s.packets[i],out.first,model,g.slice);
            project(s.packets[j],out.second,model,g.slice);
            if (a.rest_weight != b.rest_weight) {
                rb::require(s.next_id < std::numeric_limits<std::uint64_t>::max(),"Packet ID overflow.");
                auto remainder = a.rest_weight > b.rest_weight ? a : b;
                remainder.rest_weight -= scattered_weight;
                remainder.id = s.next_id++;
                s.packets.push_back(remainder);
            }
\end{lstlisting}
\end{minipage}

Projection stores specific radial and tangential momenta in the maintained
marker convention. In the inward EF branch the reciprocal expression
preserves positive canonical momentum while avoiding a subtraction of
nearly equal energies.
\par\medskip\noindent
\begin{minipage}{\linewidth}
\codesource{Specific-momentum projection}{\path{cpu/include/ev/relativistic_boltzmann/spherical_collisions.hpp}, lines 84--94}{Return the local outgoing momentum to SC or positive EF canonical variables}
\begin{lstlisting}[language=C++,showstringspaces=false]
    const double w = k.p.x/model.material_mass;
    const double transverse2 = (k.p.y*k.p.y+k.p.z*k.p.z)/(model.material_mass*model.material_mass);
    p.eta = 0.5*p.r*p.r*transverse2;
    if (slice == Slice::schwarzschild) {
        p.radial = w;
    } else {
        const double epsilon = k.t/model.material_mass;
        p.radial = w >= 0 ? epsilon+w : (1+transverse2)/(epsilon-w);
        rb::require(p.radial > 0 && std::isfinite(p.radial),"EF projection lost the future shell.");
    }
}
\end{lstlisting}
\end{minipage}



After either acceptance or rejection, a fresh exponential optical depth
starts the next proposal. Empty cells, one-packet cells and zero-area
calls advance their coordinate clocks without collision draws. The local
clock update is literal:
\par\medskip\noindent
\begin{minipage}{\linewidth}
\codesource{Renewal after every proposal}{\path{cpu/include/ev/relativistic_boltzmann/spherical_collisions.hpp}, lines 246--250}{Renew after accepted or rejected proposals and commit the local clock}
\begin{lstlisting}[language=C++,showstringspaces=false]
        s.optical_depth = -std::log(s.random.uniform());
    }
    s.coordinate_time = end;
    validate(s,model,g);
    caller = std::move(s);
\end{lstlisting}
\end{minipage}

An unequal-weight remainder receives a new identifier. Its ancestor points
to the incoming marker record, so auxiliary properties and source labels
survive replacement. The next unused identifier advances across the whole
cell traversal. The caller owns its persistence together with every cell
clock and the global transport state.

\subsection{Outgoing stress, constraints and transaction acceptance}

After all cell batches have supplied outgoing markers, a new constraint
solve reconstructs their gravitational source. Each collision batch used
the same streamed endpoint metric. The SC solve retains its positive
polar-areal gap guard. EF retains the complete interior population and
the isolated vacuum outer guard.
The reported mass difference measures the finite collision/deposition
step before the global clock is advanced.
\par\medskip\noindent
\begin{minipage}{\linewidth}
\codesource{Outgoing SC source reconstruction}{\path{cpu/include/ev/relativistic_boltzmann/coupled_evolution.hpp}, lines 161--170}{Measure the ADM residual and commit the accepted global time}
\begin{lstlisting}[language=C++,showstringspaces=false]
    out.field = ev::deposit_and_solve_weighted_einstein_vlasov_pic_field(
        s.markers,grid,evolution.shape,evolution.threads);
    rb::require(out.field.metric.min_gap > grid.parameters.gap_stop,"Postcollision SC constraint gap reached its stop.");
    // Local elastic four-momentum conservation and global ADM conservation
    // are different checks. Reconstructing the field measures the ADM change
    // associated with this frozen collision substep.
    out.collision_adm_change = out.field.adm_mass-out.streaming.endpoint_field.adm_mass;
    s.time += out.streaming.accepted_dt;
    caller = std::move(s);
    return out;
\end{lstlisting}
\end{minipage}

The EF entry point uses the same accepted-interval rule with its null rate.
It constructs radial constraint options that retain the complete interior
population and first measures a regular-centre ADM value.
It then reapplies the isolated conserved-ADM boundary condition. Keeping
both values exposes the source-feedback residual to the caller.
\par\medskip\noindent
\begin{minipage}{\linewidth}
\codesource{Outgoing EF constraints and boundary condition}{\path{cpu/include/ev/relativistic_boltzmann/coupled_evolution.hpp}, lines 235--251}{Retain every particle and expose the independently reconstructed ADM residual}
\begin{lstlisting}[language=C++,showstringspaces=false]
    // Independent regular-center reconstruction exposes any collision/sampling
    // ADM residual before the isolated fixed-ADM boundary is reapplied.
    // With no specified outer ADM value, this reconstruction integrates from
    // the regular center and measures the mass implied by the outgoing source.
    const auto regular = ef::deposit_and_solve_field(s.markers,grid,evolution.shape,evolution.threads,constraint);
    out.regular_center_reconstructed_adm = regular.metric.adm_mass;
    out.collision_adm_change = regular.metric.adm_mass-out.streaming.conserved_adm_mass;
    // The following solve instead fixes the outer ADM integration constant
    // to the streaming value. Its inner mass condition is therefore a distinct
    // boundary choice; the regular reconstruction above remains the source
    // normalization diagnostic.
    constraint.fixed_adm_mass = out.streaming.conserved_adm_mass;
    // After every cell has returned its outgoing momentum, solve the radial
    // constraints with the declared boundary choice. The clock advances by
    // the actual accepted streaming interval, including a shortened step.
    out.field = ef::deposit_and_solve_field(s.markers,grid,evolution.shape,evolution.threads,constraint);
    s.time += out.streaming.accepted_dt;
\end{lstlisting}
\end{minipage}

\subsection{Amplitude interface, source map and checks}

A supplied reciprocal scalar amplitude can determine a local differential
cross section. Write \(\mathcal M(x,s,t)\) for this amplitude, with
\(s=-(k_1+k_2)^2\) the squared centre-of-momentum energy and
\(t=-(k_1-k_3)^2\) the momentum-transfer invariant. The squares are
metric contractions. With standard quantum field theory (QFT)
external-state normalization and
\(\hbar=c=1\), elastic scattering has
\[
 \frac{\dd\Sigma}{\dd\Omega}
   =\frac{|\mathcal M(x,s,t)|^2}{64\pi^2s\,S_{\rm f}},\qquad
 \dd\Pi_{\rm QFT}=\frac{\dd^3k}{2E(2\pi)^3}.
 \tag{REB13}
\]
Here \(S_{\rm f}=2!\) for
identical final particles integrated over the full sphere. The supplied
amplitude already includes physical exchange symmetry. This final-state
factor is separate from incoming unordered-pair counting. The local
sampler uses a certified differential envelope and joint angular/rate
thinning for a bounded anisotropic callback. A chosen scalar-contact
benchmark supplies a nonconstant total cross section proportional to
\(1/s\). Curvature-dependent coefficients describe a supplied local
interaction model. Covariance alone specifies their transformation law.
The quasiparticle, local-interaction and slowly varying background
approximations determine the physical domain of this construction.

For numerical masses and momenta expressed in energy unit \(E_{\rm unit}\),
the numerical natural-unit area converts as
\(\Sigma_{\rm code}=\Sigma_{\rm natural}
[\hbar c/(E_{\rm unit}L_{\rm unit})]^2\).
The geometrized material mass uses the same physical energy through
\(m_{\chi,\rm code}=GE_\chi/(c^4L_{\rm unit})\).
The spherical coupled sampler currently takes a constant code area.
An amplitude-dependent spherical coupled restart and its physical unit
mapping require their own model contract.

All paths in the following source map share
\path{cpu/include/ev/relativistic_boltzmann/}:
\begin{description}[style=nextline]
\item[\path{relativistic_collision.hpp}]
\path{rb::pair_geometry}, \path{rb::scatter_with_direction} and
\path{rb::scatter_isotropic} provide local mass-shell kinematics and flux.
\item[\path{local_dsmc.hpp}]
\path{rb::advance_cell} supplies proper-volume weighted collisions and
mechanical local-state serialization.
\item[\path{schwarzschild_collision_adapter.hpp}]
\path{rb::sc_cell_geometry} evaluates proper volume and lapse from the
maintained field. The adapter also lifts SC markers and computes stress.
\item[\path{coordinate_geometry.hpp}]
Local metric, tetrad and shell-measure utilities support passive general
coordinate transformations at a common event.
\item[\path{invariant_amplitude.hpp}]
The model callback, differential kernel and joint angular/rate sampler
implement supplied scalar interactions with certified bounds.
\item[\path{spherical_collisions.hpp}]
\path{eb::advance_cell} implements (REB12), conservative packet splitting
and SC/EF stress with represented rest weights.
\item[\path{ef_frame.hpp}]
\path{eb::ef_frame} implements (REB6), including its horizon-regular
observer and Jacobian for the canonical momentum \(P\).
\item[\path{eb_restart.hpp}]
\path{eb::serialize_cell} and \path{eb::deserialize_cell} bind the local
clock, packets and model to the selected slice and cell geometry.
\item[\path{coupled_evolution.hpp}]
\path{eb::advance_sc} and \path{eb::advance_ef} connect maintained EV
streaming, cell collisions and Einstein-source reconstruction.
\end{description}

Finite component validation covers horizon-regular EF frames, absolute null
collision rates, local conservation, model-bound cell restart, collision-off
EV identity and small SC/EF source-feedback steps. Kinematic and amplitude
checks cover boosts, coordinate transformations and angular/rate sampling.
Each call also enforces its finite conditioning, capacity and conservation
guards. The implementation regression sources are maintained separately from
the distributed solver and analysis package.

\subsection{Spatial locality, restart and scientific domain}

The default spatial policy requires a co-located radial population. This
represents a local spherical distribution with sampled tangential momentum
azimuths. A finite radial collision cell requires explicit opt-in and
records every separated radial exchange. Distinct radii on an SC slice
are spacelike separated. Distinct radii on one EF null slice lie along a
null surface, where a two-way instantaneous exchange also needs a
spatial-locality approximation. Future mass shells and a transport CFL
condition each constrain a different part of this algorithm.

Internal Poisson times parameterize the frozen collision suboperator.
A layer-to-layer causal transport demonstrator belongs to its stated
discrete domain of dependence. Continuous event-to-transport chronology
requires causal departure times, propagation delays and self-consistent
metric control.

The public entry points commit the copied state only after the full
streaming, collision, projection and constraint sequence passes. A
conditioning, capacity, boundary or source failure leaves the caller's
markers and random state at the incoming accepted state.

Local serialization restores the complete model-bound collision cell.
A full coupled restart also needs global markers, geometry, field and
evolution controls, packet IDs and every persistent cell clock. The
current finite tests establish component conservation, rates and source
feedback. Compact-halo stability, collisional convergence, physical
escape budgets, gravothermal core collapse and trapped-surface formation
remain scientific calculations requiring their own initial data,
collisionless controls and accuracy study.
